% Corrected shared-circle synthesis (v2)
\documentclass[11pt]{article}
\usepackage{series}

\hypersetup{
  unicode=true,
  pdftitle={The Shared Circle in Modal Triplet Theory},
  pdfauthor={Peter Nero},
  pdfsubject={Exact phase transport, interpretive roles, and physical boundaries}
}

\newcommand{\MTTSharedLine}{L_{\mathrm{sh}}}
\newcommand{\MTTCLNSpace}{\mathcal H_{\mathrm{CLN}}}

\title{The Shared Circle in Modal Triplet Theory:\\
\large Exact Phase Transport, Interpretive Roles, and Physical Boundaries}
\author{Peter Nero}
\date{July 2026\\Version 2}

\begin{document}
\maketitle

\begin{abstract}
Modal Triplet Theory repeatedly uses a common circle in its bundle,
proto-spinor, finite-operator, and compactification descriptions. This paper
states what that circle currently means and what it does not yet explain. Its
rigorous role is common phase and holonomy data carried by a Hermitian line
bundle with connection. On the selected q79 finite symbol, one universal flat
line pulls back to the determinant-sign sector, tensors the three trace lanes,
and commutes with the finite projectors, complex quarter-turn, and
Reynolds-Hessian operator. This is an exact same-source phase-transport
result. It neither adds a seventh compact coordinate nor identifies compact
phase with Lorentzian time. A physical mass claim still requires a selected
Dirac or two-point operator and a dispersion or pole theorem; inertia requires
an effective response or worldline coefficient; gravity requires the
four-dimensional action and universal stress response; and time requires a
Lorentzian evolution and an explicit clock or ordering map. The proposed
unification of mass, inertia, gravity, and time is therefore retained as a
research program organized by one shared differential line, not as a current
derivation of those observables.
\end{abstract}

\section*{Version 2 Revision Note}
\begin{description}
\item[Supersedes:] The first release under the title
\emph{The Central Circle: Inertia, Mass, Gravity, and Time as Shared
Coherence Bookkeeping in Modal Triplet Theory}.
\item[Reason:] A common compact phase was identified directly with mass,
inertia, gravity, physical time, family multiplicity, and irreversible
dynamics without the required source operators or readout maps.
\item[Resolution:] The paper is rebuilt around the exact finite q79
shared-line theorem and gives a separate physical contract for each proposed
readout.
\item[Retained result:] One flat differential line supplies common connection
and holonomy data to the finite q79 lanes and Hessian square without a fitted
dimensionless coefficient.
\item[Remaining boundary:] The nonflat HYM continuation, physical mass pole,
inertial response, four-dimensional stress coupling, clock map, and
branch-selection law remain open.
\end{description}

\tableofcontents

\section{The answer first}

The phrase ``shared circle'' has been used in several ways across the MTT
corpus. The current reconciled meaning is:
\[
\boxed{\text{one source of compact \(U(1)\) phase and holonomy data}}
\]
represented by a Hermitian line bundle and its unitary connection.

This statement is global enough to compare transport around loops. It is
stronger than noticing several abstract circles with the same dimension, but
weaker than saying that all physical sectors are generated by one coordinate.
The circle is counted once. It is not multiplied through the Lens, Nil, and
q79 factors, and it is not an extra product dimension.

The exact current theorem belongs to the corrected MTT Foundation. It proves a
finite shared-line and Hessian intertwiner on the selected q79 flat
root-stack symbol. The present paper explains that result and asks how far it
can be transported toward familiar physical concepts.

\subsection{Four levels that must not be collapsed}

\begin{center}
\begin{tabular}{p{0.25\linewidth}p{0.30\linewidth}p{0.34\linewidth}}
\hline
Level & Mathematical object & Current status\\
\hline
Phase fiber & A copy of \(U(1)\) in each fiber & Standard construction\\
Differential line & Hermitian line plus connection and holonomy & Exact on
the finite q79 symbol\\
Continuum coupling & Line coupled to HYM, Dirac, or four-dimensional fields &
Partly specified; physical intertwiner open\\
Observable readout & Mass, inertia, gravity, clock time, or family data &
Requires a separate theorem for each readout\\
\hline
\end{tabular}
\end{center}

The older manuscript moved directly from the first row to the fourth. The
finite shared-line result closes the second row and gives a useful source for
the third. It does not close the observable rows automatically.

\section{Circle fibers, line bundles, and connections}

A bare circle \(S^1\) is a one-dimensional compact manifold. In gauge and
bundle geometry, the more useful object is a principal \(U(1)\) bundle or its
associated Hermitian line bundle. Locally it resembles a product
\(U\times U(1)\), but globally its transition functions can carry nontrivial
topology and holonomy.

A unitary connection \(\nabla\) tells us how phase is transported. Its
curvature records infinitesimal transport around small loops, while holonomy
records transport around finite loops. Two line bundles with isomorphic
fibers need not have the same topology, connection, or holonomy. Thus the
claim that different sectors use ``the same circle'' is meaningful only after
the source line and the connection-preserving pullback maps are named.

\subsection{What a shared source guarantees}

Suppose a Hermitian line with connection \((L,\nabla)\) is defined over a
classifying base and sector maps \(f_i:X_i\to B\) are supplied. The sector
lines are the pullbacks
\[
 (L_i,\nabla_i)=f_i^*(L,\nabla).
\]
Their curvatures and loop holonomies are then pullbacks of one differential
object. This gives common phase provenance. It does not make the spaces
\(X_i\) identical or make every operator acting on them the same.

An operator comparison needs another square: the line pullbacks must tensor
the relevant state bundles, and the sector operators must intertwine the
pullback maps. That extra operator square is exactly what the current q79
finite result supplies.

\section{The exact q79 finite shared-line result}

The selected finite branch uses a degree-three sheet system with monodromy in
\(S_3\). Its abelianization has one nontrivial sign character. The corrected
Foundation maps that sign into the order-two element of a universal flat
\(\mathbb Z_{64}\) differential line. In the Foundation's notation,
\[
 h_{S_3}(\sigma)=32\,\epsilon(\sigma),
\]
where \(\epsilon(\sigma)\) distinguishes even and odd permutations. For both
admissible odd roots, the resulting character agrees with the determinant
sign. This produces a parallel isomorphism, not merely a fiberwise one:
connections and every loop holonomy agree.

The same pulled-back line \(\MTTSharedLine\) then tensors the three trace lanes,
\[
 \MTTCLNSpace
 =
 \Lsh\otimes
 \left(
  \mathcal O\oplus\mathcal A_0\oplus\mathcal A
 \right),
\]
whose ranks are \(1\), \(2\), and \(3\). Scalar phase transport commutes with:
\begin{enumerate}
\item the trace, trace-free, and full projectors;
\item the root-plane quarter-turn \(J_{DE}\);
\item the finite Haar or Reynolds projector; and
\item the normalized finite Hessian
  \(H_{\mathrm{fin}}=\kappa_{\mathrm{fin}}(I-P_{\mathrm{Haar}})\).
\end{enumerate}

On the selected two-copy sheet/edge representation, the finite calculation
gives
\[
 \operatorname{rank}P_{\mathrm{Haar}}=2,
 \qquad
 \operatorname{spec}(H_{\mathrm{fin}}/\kappa_{\mathrm{fin}})
 =\{0^{\times2},1^{\times4}\},
 \qquad
 H_{\mathrm{TT}}=\kappa_{\mathrm{fin}}I_2.
\]
Tensoring by \(\MTTSharedLine\) leaves these relations unchanged. No dimensionless fit
is introduced to make the phase and Hessian data compatible.

\subsection{Why this is a real advance}

Earlier papers often said that several sectors ``reuse a circle.'' That could
have meant merely that each sector had some isomorphic copy of \(U(1)\). The
finite q79 theorem instead identifies one differential source, its
connection, its holonomy, and its commuting action on the relevant finite
operators. It is a same-source result.

The theorem is exact at flat differential-character and finite-symbol tier.
The physical HYM bundle has nonzero characteristic curvature, so its
connection cannot literally equal this flat line connection. The next
continuum result must be an intertwining or spectral-symbol functor that
respects the different curvature data.

\section{Lens, Nil, and the shared-circle language}

Lens and Nil geometries remain useful, but not as literal nested containers.
Both can arise from the Boothby--Wang or prequantization circle-bundle
construction over different bases. Over a two-sphere, an integral circle
bundle can give a lens space; over a two-torus, it can give a Heisenberg
nilmanifold. Their connection curvatures and total-space topologies differ.

The reconciled picture is therefore:
\[
\begin{array}{c}
\text{flat q79 line: finite sign/phase transport},\\[2mm]
\text{Lens and Nil bundles: curved circle-bundle realizations},\\[2mm]
\text{physical HYM bundle: nonflat holomorphic gauge data}.
\end{array}
\]
These objects share a construction language through differential
\(U(1)\)-bundle geometry. They are not one manifold, one connection, or one
extra coordinate. Literal \(S^1\times\mathrm{Lens}\times\mathrm{Nil}\) and
literal circle-inside-Lens-inside-Nil nesting are not used as the physical
q79 compactification.

\section{Mass: what a circle can and cannot supply}

\subsection{The physical definition comes first}

In relativistic quantum theory, a stable particle mass is read from the
energy-momentum representation or, in a field description, from the location
of the appropriate propagator pole. For unstable states one uses a complex
pole or another explicitly declared scheme. A running Lagrangian parameter,
a Hessian eigenvalue, and a physical pole mass are not automatically the same
number.

Compact phase data can influence masses. Boundary conditions or holonomy can
shift the spectrum of a Dirac, Laplace, or covariant derivative operator.
Dimensional reduction can then turn an internal eigenvalue into a
four-dimensional mass parameter. But the complete chain must be displayed:
\[
\begin{gathered}
\text{line and connection}
\longrightarrow
\text{coupled operator}
\longrightarrow
\text{eigenvalue or self-energy}\\
\longrightarrow
\text{dispersion relation or pole}.
\end{gathered}
\]

The finite shared-line theorem closes only the first arrow and its
compatibility with a finite Hessian. It does not identify
\(\kappa_{\mathrm{fin}}\) with an observed particle mass.

\subsection{Relation to the Standard-Model carrier}

The current MTT Standard-Model program has a finite \(27\times27\) carrier and
declared Yukawa, threshold, and profile/value tiers. Those results can
constrain an internal mass matrix. To promote a row to a physical mass one
still needs the normalized four-dimensional kinetic term, renormalization
scheme and scale, self-energy or matching transport, and the pole or
dispersion calculation.

The shared line is a promising common source for phases and selection rules.
The statement that it numerically generates all masses or Yukawa magnitudes
is not presently proved.

\section{Inertia: response is not phase alone}

Inertia is the response of a physical state to changes in motion. For a
relativistic point-particle model, the action
\[
 S_{\mathrm{pp}}=-m\int ds
\]
has a low-velocity expansion whose coefficient \(m\) produces the familiar
\(\mathbf F=m\mathbf a\) relation. Deriving that relation from the assumed
action is standard. It does not derive \(m\) from a circle.

An MTT inertia theorem would need to start from the selected field or
effective action, identify a localized state, and compute its response to an
external force or to a change in four-momentum. The result would then have to
agree with the mass extracted from the pole or dispersion relation in the
same normalization.

This gives a sharp equality test:
\[
 m_{\mathrm{response}}
 \stackrel{?}{=}
 m_{\mathrm{pole}}
 \stackrel{?}{=}
 m_{\mathrm{gravity}} .
\]
The equal signs are physical theorems, not consequences of using one phase
line in all three calculations.

\section{Gravity: a common line is not the Einstein equation}

The current MTT gravity program conditionally recovers the
four-dimensional Einstein/TEGR tensor structure from a selected higher
dimensional branch after truncation, Lorentzian, normalization, and
stress-response hypotheses are supplied. That result is the proper source
for gravitational dynamics.

The shared line may participate in the compactification, spin structure,
phase transport, or anomaly data entering that reduction. It can therefore
help enforce common provenance across matter and gravity sectors. It does not
by itself imply:
\begin{enumerate}
\item the Einstein--Hilbert or TEGR action;
\item universal coupling to the full stress tensor;
\item equality of inertial and gravitational mass;
\item the Newton normalization; or
\item horizon entropy and information recovery.
\end{enumerate}

Each item requires the four-dimensional action or response functional. The
shared-circle interpretation of gravity as ``coherence bookkeeping'' remains
a useful conceptual proposal only if the selected action later realizes it.

\section{Time: phase, order, and clocks are different objects}

Four parameters have appeared near the circle discussion:
\begin{enumerate}
\item the compact phase angle \(\theta\);
\item a stabilization parameter \(\tau\);
\item physical Lorentzian time \(t\); and
\item a renormalization scale \(\mu\).
\end{enumerate}
They are not interchangeable. The phase is periodic. Stabilization labels a
mathematical relaxation. Physical time belongs to a Lorentzian initial-value
problem. Renormalization compares descriptions at different scales.

\subsection{How a circle may participate in a clock}

A periodic degree of freedom can serve as part of a clock if a physical
Hamiltonian, clock observable, calibration, and method of counting or
unwrapping cycles are supplied. Even then, a periodic clock generally gives
only local or cycle-relative time unless additional nonperiodic information
is available. The circle phase is therefore capable of synchronizing sectors,
but it is not itself a global noncompact time coordinate
\cite{QuantumClocks,PeriodicClocks}.

\subsection{Projection does not automatically create an arrow}

A many-to-one projection loses upper distinctions, but the induced lower
dynamics depends on how the upper evolution acts on projection fibers. The
correct quotient theorem is:
\begin{quote}
an autonomous lower map exists exactly when upper evolution preserves
projection fibers; it is invertible exactly when the fiber equivalence is
preserved in both directions.
\end{quote}
Thus a noninjective projection can coexist with reversible lower dynamics.
The old claim that the absence of a global right inverse creates physical
time is withdrawn.

An arrow of time requires a selected branch or state, Lorentzian evolution,
and a monotone physical record, entropy, or clock relation. The two
complex-conjugate finite phase lifts can remain mathematically admissible
without both being well-behaved physical universes. Selecting the observed
orientation remains a separate problem.

\section{Families and phase selection}

The q79 construction has a degree-three sheet system, and its finite carrier
uses the \(1+2+3\) rank flag. The shared line carries the abelianized sign
phase of the sheet monodromy and participates in the finite SpinC/Fourier
return. These are exact finite statements.

Three fermion families do not follow from the existence of an arbitrary
circle. Their proposed identification uses the selected degree-three cover,
monodromy, bundle representations, and the Standard-Model operator map. A
circle can support many characters and windings; it does not uniquely choose
three. The family claim must therefore be sourced by the full q79 cover and
operator construction, with the shared line supplying common phase transport
rather than the family count by itself.

\section{Current status of the unification proposal}

\begin{center}
\begin{tabular}{p{0.25\linewidth}p{0.18\linewidth}p{0.46\linewidth}}
\hline
Claim & Status & Required evidence\\
\hline
One finite q79 shared differential line & Exact & Foundation finite
shared-line theorem\\
Compatibility with finite projectors and Hessian & Exact & Commuting tensor
action on selected symbol\\
Continuation to physical HYM geometry & Open & Connection-aware continuum
intertwiner with curvature respected\\
Circle determines particle mass & Open & Coupled operator plus pole or
dispersion theorem\\
Circle determines inertia & Open & Localized-state response calculation\\
Circle enforces equivalence principle & Open & Same-source inertial and
gravitational response equality\\
Circle is physical time & Rejected & Compact phase and noncompact time are
different objects\\
Circle participates in a clock & Conditional & Hamiltonian, observable,
calibration, and cycle handling\\
Circle alone gives three families & Rejected & Degree-three cover and
operator map are essential\\
Shared line organizes all four readouts & Research proposal & One
connection-preserving source through all required operators\\
\hline
\end{tabular}
\end{center}

\section{The shortest route forward}

The promising idea is not that one circle should be renamed as four physical
quantities. It is that one differential line may feed four independently
defined operators. A decisive same-source program would construct:
\begin{enumerate}
\item a continuum map from the flat q79 line to the physical
  SpinC/HYM/Dirac data, preserving the declared connection relations;
\item a normalized four-dimensional Dirac or two-point operator and a
  certified physical mass extraction;
\item an inertial response functional for the same localized state;
\item the four-dimensional gravitational response to that state's stress
  tensor; and
\item a physical clock observable whose phase readout is related to
  Lorentzian time on a declared domain.
\end{enumerate}

Agreement among these readouts would be a nontrivial unification because the
same source would survive several different maps. Disagreement would locate
the exact arrow where the shared-circle interpretation fails.

\section{Discussion}

The finite shared-line theorem is stronger than an analogy and weaker than a
theory of mass or time. It says that the q79 finite lanes and Hessian square
really do share one phase-transport source. That is valuable: phases,
determinant signs, and finite operators no longer need unrelated circle
copies.

The physical ambitions of the first manuscript remain interesting. A common
line could help correlate flavor phases, spinorial closure, compactification,
and effective response. But physics is attached to operators and observables,
not to suggestive names. Mass must be a mass in a dispersion law or
propagator; inertia must be a response; gravity must follow from a varied
action; and time must be read by a physical clock within a Lorentzian
dynamics.

\section{Conclusion}

The shared circle survives the corpus reconciliation in a precise form. It is
one Hermitian differential line carrying compact phase and holonomy, counted
once, with an exact finite q79 intertwining action. It is not an extra
dimension, literal compact time, or a direct derivation of mass, inertia,
gravity, and families.

The unifying research program is now testable: transport the same line into
the continuum operators that independently define each observable and check
whether their normalized readouts agree. That is the correct way to turn
shared bookkeeping into shared physics.

% BEGIN MTT MANAGED COMPUTATIONAL EVIDENCE
\section*{Computational Evidence and Reproducibility}
The shared-circle paper uses the selected q=79, finite-matrix, HYM, anomaly, threshold, and electroweak rows as concrete instances of common phase or holonomy data. These packets do not identify the compact circle with physical time or derive inertia, mass, or gravity. The profile rows are context and the strict no-knob upgrade remains open.

The referenced rows are frozen to the curated results repository state identified below. Hashes are grouped in eight-character blocks for line breaking.
\begin{quote}\small
\textbf{Repository:} \url{https://github.com/PeterNero/mtt-results-repro}\\
\textbf{Commit:} \texttt{31247ebb 5c22f3fb b5443024 365433c6 ee0bff4a}\\
\textbf{Manifest:} \path{release/result_manifest.json}\\
\textbf{Manifest SHA-256:}\\
\texttt{fb399689 60b00584 631dbf53 1a708e18}\\
\texttt{ef928d6b 6d935119 c185d7f6 32b1e7cd}
\end{quote}

Tier labels are quoted verbatim from that manifest. A row used directly supports only the specific computational statement identified above; a corpus-state cross-check does not prove this paper's local theorems; and an open row is evidence of an unresolved obligation, never of closure.

\par\begingroup\dimen0=\pagegoal\advance\dimen0 by -\pagetotal\ifdim\dimen0<7\baselineskip\newpage\fi\endgroup
\paragraph{Rows used directly in this paper.}
\begin{itemize}
\setlength{\itemsep}{0.25em}
\item \path{A01/direct_k_higgs_row} (\emph{derived exact}).\par Promoted direct K\_threshold.Omega\_H.lambda row.
\item \path{A22/e6_qpsi_qcd_anomaly} (\emph{derived exact}).\par E6 Qpsi matter/exotic QCD anomaly cancellation audit.
\item \path{A19/hym_wiener_contraction} (\emph{numeric certified}).\par Weighted-theta Fourier-tail and Wiener contraction certificate.
\item \path{A11/q79_exact_audit} (\emph{derived exact}).\par Executable q=79 exact-branch audit.
\item \path{A11/q79_exact_theorem} (\emph{derived exact}).\par CRT q=79 theorem on the selected exact branch.
\item \path{A01/qutrit_weyl_27_matrix} (\emph{derived exact}).\par Sparse 27x27 qutrit-Weyl left-action realization.
\item \path{A01/strict_pew_row} (\emph{derived exact}).\par Promoted P\_EW source row at the declared one-shared-primitive standard.
\end{itemize}

\par\begingroup\dimen0=\pagegoal\advance\dimen0 by -\pagetotal\ifdim\dimen0<7\baselineskip\newpage\fi\endgroup
\paragraph{Corpus-state cross-checks.}
\begin{itemize}
\setlength{\itemsep}{0.25em}
\item \path{A01/charged_yukawa_higgs_profile} (\emph{profile replay}).\par Versioned Yu, Yd, Ye and lambda\_H profile packet.
\item \path{A14/ckm_prediction_profile} (\emph{numeric certified}).\par Three selected CKM profile rows and uncertainty comparison.
\item \path{A01/current_global_lock} (\emph{profile replay}).\par Current non-looping global status and source-certificate map.
\item \path{A04/final_12_of_12_audit} (\emph{profile replay}).\par Twelve-obligation embedded renormalized-SM equivalence audit.
\item \path{A40/neutral_two_primitive_profile} (\emph{profile replay}).\par Measured two-splitting neutrino profile, masses and Dirac Yukawa rows.
\item \path{A02/precision_15_source_transport} (\emph{profile replay}).\par Fifteen measured source coordinates, Jacobian and covariance transport.
\item \path{A02/precision_8x8_workspace} (\emph{profile replay}).\par Eight-coordinate SMDR output with positive-definite 8x8 covariance.
\end{itemize}

\par\begingroup\dimen0=\pagegoal\advance\dimen0 by -\pagetotal\ifdim\dimen0<7\baselineskip\newpage\fi\endgroup
\paragraph{Open boundary (not evidence of closure).}
\begin{itemize}
\setlength{\itemsep}{0.25em}
\item \path{A05/strict_upgrade_ledger} (\emph{open}).\par Current 2/9 strict no-knob upgrade ledger.
\end{itemize}

No imported row changes theorem ownership or promotes a neighboring claim: all local statements retain their stated hypotheses, domains, and limitations.
% END MTT MANAGED COMPUTATIONAL EVIDENCE

\begin{thebibliography}{9}

\bibitem{MTTFoundation}
P.~Nero,
\emph{Modal Triplet Theory: Foundations},
current revised MTT paper corpus, 2026.

\bibitem{MTTProgramC}
P.~Nero,
\emph{The Modal Triplet Theory Program C: Realizing the Modal Carrier},
current revised MTT paper corpus, 2026.

\bibitem{FuYau}
J.-X.~Fu and S.-T.~Yau,
\emph{The theory of superstring with flux on non-Kahler manifolds and the
complex Monge--Ampere equation},
Journal of Differential Geometry \textbf{78} (2008), 369--428.

\bibitem{PoleMass}
A.~K.~Das, R.~R.~Francisco, and J.~Frenkel,
\emph{On the gauge independence of the fermion pole mass},
Physical Review D \textbf{88}, 085012 (2013).

\bibitem{QuantumClocks}
M.~B.~Altaie, D.~Hodgson, and A.~Beige,
\emph{Time and quantum clocks: a review of recent developments},
Frontiers in Physics \textbf{10}, 897305 (2022).

\bibitem{PeriodicClocks}
L.~Chataignier, P.~A.~Hoehn, M.~P.~E.~Lock, and F.~M.~Mele,
\emph{Relational dynamics with periodic clocks},
arXiv:2409.06479.

\end{thebibliography}

\end{document}
